\clearpage
\section{Phân phối chuẩn nhiều chiều: tâm, hình dạng và khoảng cách}
\label{sec:normal}

\subsection{Hai tham số đủ mô tả mô hình}
Trong một chiều, $\N(\mu,\sigma^2)$ có tâm $\mu$ và độ phân tán
$\sigma^2$. Trong $p$ chiều, vai trò tương ứng thuộc về vectơ $\vmu$ và
ma trận $\Sigma$:
\[
 \mathbf X\sim\N_p(\vmu,\Sigma).
\]
Vectơ $\vmu$ đặt tâm của đám mây dữ liệu. Ma trận $\Sigma$ vừa cho độ phân
tán của từng biến vừa cho cách chúng cùng biến thiên. Vì vậy biết các trung
bình và các phương sai riêng vẫn chưa đủ: các hiệp phương sai cũng cần được xác định.

\begin{definition}[Chuẩn nhiều chiều không suy biến]
Với $\vmu\in\R^p$ và $\Sigma$ đối xứng xác định dương, mật độ tại
$\mathbf x\in\R^p$ là
\begin{equation}
 f(\mathbf x)=\frac{1}{(2\pi)^{p/2}|\Sigma|^{1/2}}
 \exp\!\left\{-\frac12(\mathbf x-\vmu)'\Sigma^{-1}
 (\mathbf x-\vmu)\right\}.
 \label{eq:density}
\end{equation}
\end{definition}
\nguon{Chương 4, §4.2, công thức (4-4), tr.~150}

\hieu{Mật độ lớn nhất ở tâm. Càng đi xa tâm theo thước đo có tính đến
phương sai và tương quan, mật độ càng giảm. Hệ số ở trước hàm mũ làm tổng
xác suất trên toàn không gian bằng 1.}

Giá trị mật độ không phải xác suất của một điểm. Với biến liên tục,
$\Prob(\mathbf X=\mathbf x)=0$; xác suất của một vùng $B$ là
$\int_B f(\mathbf x)\,d\mathbf x$. Trong hai chiều, đây là thể tích dưới
bề mặt mật độ, trên vùng $B$.

\subsection{Đọc mật độ mà không cần thuộc lòng}
Hãy tách công thức thành ba thành phần:
\begin{enumerate}
 \item $\mathbf x-\vmu$: độ lệch khỏi tâm.
 \item $\Sigma^{-1}$: hiệu chỉnh độ lệch theo độ phân tán và tương quan.
 \item $\exp(-D^2/2)$: làm mật độ giảm khi độ lệch chuẩn hóa $D^2$ tăng.
\end{enumerate}
Định thức $|\Sigma|$ thường được gọi là \emph{phương sai tổng quát}.
Nếu độ phân tán theo các hướng tăng, cùng một lượng xác suất trải ra trên
thể tích lớn hơn và đỉnh mật độ thấp hơn. Tuy nhiên định thức không cho biết
riêng từng hướng: hai ma trận khác nhau có thể cùng định thức.

Nếu $p=1$, $|\Sigma|=\sigma^2$ và công thức trở lại
\[
 f(x)=\frac1{\sqrt{2\pi}\sigma}
 \exp\left[-\frac12\left(\frac{x-\mu}{\sigma}\right)^2\right].
\]
Đây là một kiểm tra hữu ích: công thức nhiều chiều phải thu về công thức
quen thuộc khi chỉ còn một biến.

\subsection{Khoảng cách Mahalanobis}
\begin{definition}
Bình phương khoảng cách Mahalanobis từ $\mathbf x$ đến $\vmu$ là
\begin{equation}
 D^2(\mathbf x,\vmu)=(\mathbf x-\vmu)'\Sigma^{-1}(\mathbf x-\vmu).
 \label{eq:mahalanobis}
\end{equation}
\end{definition}
Nếu $\Sigma=\diag(\sigma_1^2,\ldots,\sigma_p^2)$, thì
\[
 D^2=\sum_{i=1}^p\left(\frac{x_i-\mu_i}{\sigma_i}\right)^2.
\]
Ta nhận được tổng các bình phương độ lệch chuẩn hóa. Nếu các biến tương
quan, phép chuẩn hóa còn phải tính đến hướng biến thiên chung.

\begin{example}[Cùng khoảng cách thường, khác mức độ bất thường]
Lấy $\vmu=(0,0)'$ và $\Sigma=\begin{pmatrix}4&3\\3&4\end{pmatrix}$.
Hai điểm $\mathbf u=(2,2)'$ và $\mathbf v=(2,-2)'$ đều có khoảng cách
Euclid $\sqrt8$ đến tâm. Nhưng
\[
 D^2(\mathbf u,0)=\frac17(2,2)
 \begin{pmatrix}4&-3\\-3&4\end{pmatrix}\begin{pmatrix}2\\2\end{pmatrix}
 =\frac87,
\]
\[
 D^2(\mathbf v,0)=\frac17(2,-2)
 \begin{pmatrix}4&-3\\-3&4\end{pmatrix}\begin{pmatrix}2\\-2\end{pmatrix}
 =8.
\]
Tương quan ở đây là $3/4$. Hai biến tăng cùng nhau là tình huống thường gặp;
một biến tăng còn biến kia giảm là hướng ít biến thiên hơn.
Vì vậy $\mathbf v$ bất thường hơn $\mathbf u$ dù khoảng cách Euclid bằng nhau.
\end{example}

\begin{figure}[htbp]
\centering
\begin{tikzpicture}[scale=.75,>=Stealth]
 \draw[->] (-3.8,0)--(4,0) node[right]{$x_1$};
 \draw[->] (0,-3.4)--(0,3.6) node[above]{$x_2$};
 \draw[rotate=45] (0,0) ellipse (3.25 and 1.23);
 \draw[dashed,rotate=45] (0,0) ellipse (2.1 and .794);
 \draw[dotted] (-3.2,-3.2)--(3.2,3.2);
 \fill (0,0) circle (2pt) node[below left]{$\vmu$};
 \fill (2,2) circle (2.4pt) node[above left]{$\mathbf u$};
 \fill (2,-2) circle (2.4pt) node[below right]{$\mathbf v$};
\end{tikzpicture}
\caption{Đường đồng mức của chuẩn hai chiều với tương quan dương.
Hướng $x_1=x_2$ biến thiên nhiều hơn hướng $x_1=-x_2$. Các đường được vẽ
để minh họa hình dạng, không gán mức xác suất cụ thể.}
\label{fig:correlation}
\end{figure}

\subsection{Tại sao đường đồng mức là elip?}
Mật độ giữ nguyên khi $D^2$ giữ nguyên. Do đó các điểm có cùng mật độ nằm trên
\[
 (\mathbf x-\vmu)'\Sigma^{-1}(\mathbf x-\vmu)=c^2.
\]
Trong hai chiều đây là đường elip, trong nhiều chiều là mặt elipsoid.
Gọi $\lambda_i$ là trị riêng của $\Sigma$ và $\mathbf e_i$ là vectơ riêng
đơn vị tương ứng. Trục thứ $i$ đi theo hướng $\mathbf e_i$, với bán trục
$c\sqrt{\lambda_i}$. Hướng có trị riêng lớn là hướng đám mây trải rộng.
\nguon{Chương 4, §4.2, công thức (4-7) và Result 4.1 về trị riêng của nghịch đảo, tr.~153}

Bạn có thể hiểu trị riêng ở đây như ``phương sai sau khi xoay hệ trục đúng
hướng''. Không cần tính trị riêng để làm mọi bài, nhưng ý nghĩa này giúp đọc hình.
Với ma trận trong ví dụ trên, hai hướng $(1,1)'/\sqrt2$ và $(1,-1)'/\sqrt2$
có phương sai lần lượt 7 và 1. Tỉ số bán trục là $\sqrt7$.

\begin{figure}[htbp]
\centering
\begin{tikzpicture}[>=Stealth,scale=.7]
 \begin{scope}[xshift=-3.2cm]
 \draw[->] (-2.5,0)--(2.5,0) node[right]{$x_1$};
 \draw[->] (0,-2.2)--(0,2.2) node[above]{$x_2$};
 \draw (0,0) circle (1.8);
 \draw[dashed] (0,0) circle (1);
 \node[align=center,text width=3.6cm] at (0,-3) {$\rho=0$\\phương sai bằng nhau};
 \end{scope}
 \begin{scope}[xshift=3.4cm]
 \draw[->] (-2.5,0)--(2.5,0) node[right]{$x_1$};
 \draw[->] (0,-2.2)--(0,2.2) node[above]{$x_2$};
 \draw[rotate=-45] (0,0) ellipse (2.3 and .8);
 \draw[dashed,rotate=-45] (0,0) ellipse (1.3 and .452);
 \node[align=center,text width=3.6cm] at (0,-3) {$\rho<0$\\elip nghiêng xuống};
 \end{scope}
\end{tikzpicture}
\caption{Tương quan quyết định hướng nghiêng; phương sai và tương quan cùng
quyết định độ dài các trục. Nếu tương quan bằng 0 nhưng phương sai khác nhau,
đường đồng mức vẫn là elip, không nhất thiết là đường tròn.}
\end{figure}

\subsection{Từ khoảng cách đến xác suất: phân phối chi-bình phương}
\begin{result}
Nếu $\mathbf X\sim\N_p(\vmu,\Sigma)$ và $\Sigma$ xác định dương, thì
\begin{equation}
 (\mathbf X-\vmu)'\Sigma^{-1}(\mathbf X-\vmu)\sim\chi^2_p.
 \label{eq:chisqdistance}
\end{equation}
Vì vậy miền
\[
 B=\{\mathbf x:D^2(\mathbf x,\vmu)\le\chi^2_{p;1-\alpha}\}
\]
chứa một quan sát ngẫu nhiên với xác suất $1-\alpha$.
\end{result}
\nguon{Chương 4, §4.2, Result 4.7(a,b), tr.~163--164}

\hieu{Sau khi chuẩn hóa đúng cả độ phân tán lẫn tương quan, ta có $p$
biến chuẩn độc lập với trung bình 0 và phương sai 1. Tổng bình phương của chúng
chính là biến chi-bình phương với $p$ bậc tự do.}
Đây là ý tưởng của phép \emph{làm trắng}; chứng minh được đặt ở
Phụ lục~\ref{app:whitening}.

\begin{example}[Một vùng xác suất 95\%]
Trong hai chiều, $\chi^2_{2;0.95}\approx5.9915$.
Điểm $\mathbf u$ ở ví dụ trên có $D^2=1.1429<5.9915$, còn
$\mathbf v$ có $D^2=8>5.9915$. Điểm thứ nhất nằm trong elip xác suất 95\%,
điểm thứ hai nằm ngoài.
Một điểm nằm ngoài không phải điều bất khả: theo mô hình, khoảng 5\% các
quan sát nằm ngoài vùng này.
\end{example}

\luuy{Đây là vùng cho \emph{một quan sát} khi biết $\vmu,\Sigma$.
Miền tin cậy cho \emph{trung bình tổng thể} khi ước lượng từ mẫu là một bài
toán khác, có $n$, $S$ và ngưỡng Hotelling. Không đổi tên vùng xác suất này
thành miền tin cậy cho $\vmu$.}

\kiemtra{Trong hai chiều, chọn ngưỡng 1 thay vì $\chi^2_{2;0.95}$ có cho
vùng 95\% không? Không: số 1 chỉ ấn định khoảng cách, chưa gắn đúng xác suất.}
